\section{Introduction}

Using LLMs as automated judges has become standard practice for evaluating model outputs~\cite{zheng2023judging,tan2024judgebench,li2024arenahard}. On text-centric tasks - dialogue, summarization, instruction following - judge reliability is well-characterized, with documented biases and known failure modes~\cite{zheng2023judging,wang2024fair}. As LLMs are now widely deployed as autonomous agents that invoke tools and orchestrate multi-step workflows, the judge paradigm has extended naturally to evaluate agentic tool-calling~\cite{qin2023toolllm,guo2024stabletoolbench,mcpagentbench2025}. This extension, however, has proceeded without a basic calibration check: \emph{how reliable are LLM judges in this structured setting?}

Agentic tool-calling differs from text evaluation in ways that matter for judge reliability. Correctness here is not a matter of fluency or preference - it requires selecting the right tools from a typed schema, supplying well-formed arguments, ordering calls to respect execution dependencies, and covering all parts of the user's intent. A plan can fail in four orthogonal ways, and these failure modes do not correlate cleanly. A judge reliable at detecting wrong tool selection may be blind to argument-key errors. A judge calibrated on simple sequential tasks may fail on complex fan-in or diamond workflows where parallel branches must converge. And when no ground-truth execution trace is available - the common deployment scenario - the judge must reconstruct correctness from the query and tool schemas alone, a fundamentally harder inference problem. Existing benchmarks that deploy LLM judges for tool-calling report only aggregate pass-rate agreement and vary none of these dimensions~\cite{qin2023toolllm,guo2024stabletoolbench,mcpagentbench2025}. The practical consequence is that practitioners deploying judge-based evaluation pipelines have no principled basis for choosing a judge, setting its configuration, or interpreting its outputs.

We introduce \textbf{AgentJudgeBench} to close this gap. The benchmark comprises 3,808 BFCL-style records~\cite{patil2025bfcl} spanning six DAG topologies at three controlled difficulty tiers, each with programmatically verified ground-truth traces. Five generators - four open-weight (3B-70B) and one frontier (GPT-5.4) - produce tool-calling outputs that are scored by six LLM judges - ranging from 20B open-weights to frontier closed systems - under paired with-GT and without-GT conditions across four structural metrics (tool selection, parameter structure, sequence accuracy, query coverage), yielding $321{,}648$ completed LLM-judge evaluations (Appendix~\ref{app:eval-stats}). We investigate seven research questions: which metrics and topologies are hardest (RQ1), how much judges agree with each other (RQ2), why no-GT alignment converges on hard queries (RQ3), and how sensitive alignment is to judge temperature (RQ4), model scale (RQ5), prompt format (RQ6), and chain-of-thought reasoning (RQ7).

The results are non-obvious. Judge alignment degrades monotonically with query difficulty - $1.5\times$ faster without ground truth than with it - and all six judges converge to a $77$-$82\%$ ceiling on hard queries without a reference, regardless of model capacity. Ground-truth exposure is not universally beneficial: two frontier judges (GPT-5.4, Gemini-2.5-Pro) are \emph{less} aligned when shown the reference, consistent with over-anchoring to the provided answer rather than exercising independent judgement. Chain-of-thought reasoning in the judge adds at most $0.3$\,pp across 24 paired comparisons, but structured per-metric prompt rubrics add $+4.8$-$+6.5$\,pp over free-form alternatives - a first-order effect. Inter-judge agreement is moderate ($\kappa \approx 0.42$ with-GT) with systematic, capacity-correlated patterns that are not visible in aggregate alignment scores.

\noindent\textbf{Contributions.} \textbf{(1)}~A dataset of 3,808 records across six DAG topologies and three difficulty tiers with programmatically verified ground-truth traces, released publicly. \textbf{(2)}~A four-metric evaluation framework with a paired with-GT/without-GT protocol and bootstrap confidence intervals. \textbf{(3)}~A systematic seven-RQ empirical study producing actionable guidance: with ground truth available, QwQ-32B dominates; without it, frontier judges lead narrowly but the convergence ceiling limits the practical difference.
